% Seminar 3 - Pagina 5. Fragment pentru
% % Seminar 3 - Pagina 5. Fragment pentru
% % Seminar 3 - Pagina 5. Fragment pentru
% % Seminar 3 - Pagina 5. Fragment pentru
% \input{seminars/pages/seminar3-pagina5-ro}; fără preambul sau pagină nouă.
% NOTE EDITORIALE:
% - Primele rânduri continuă proprietatea de universalitate de pe pagina 4.
%   Termenul din al doilea criteriu este suprascris și poate fi „inj.” ori
%   „bij.”. Am adoptat provizoriu „injectivă”; dacă scria „bijectivă”,
%   ar trebui inclusă și condiția de surjectivitate.
% - În exemplul numerelor raționale apare Z x Z, fără excluderea
%   numitorului zero. Înmulțirea în cruce nu dă o relație de echivalență
%   pe întregul Z x Z; nu am adăugat tacit restricția necesară.
% - Șirurile Cauchy sunt interpretate ca șiruri cu termeni în Q,
%   potrivit notației compacte din manuscris.
% - Schița proiecției din exemplul (i) a fost refăcută ca diagramă.
% - Însemnările foarte palide din marginea dreaptă și cercul hașurat
%   nu au un rol matematic sigur și nu sunt incluse în textul principal.
% - Identificarea R/~ cu R de la (iv) este abstractă, nu canonică.

\subsection*{Seminar 3 - Pagina 5}
\begingroup
\setlength{\abovedisplayskip}{5pt}
\setlength{\belowdisplayskip}{5pt}
\setlength{\abovedisplayshortskip}{3pt}
\setlength{\belowdisplayshortskip}{3pt}

\noindent\textbf{Proprietatea de universalitate (continuare).}
Pentru aplicația $\overline f$ de pe pagina anterioară,
\[
\begin{aligned}
 \overline f\text{ este surjectivă}
   &\iff f\text{ este surjectivă},\\
 \overline f\text{ este injectivă}
   &\iff {\sim}\;=R_f.
\end{aligned}
\]
Mai mult, pentru un sistem de reprezentanți $(x_i)_{i\in I}$,
\[
 M=\bigcup_{i\in I}\widehat{x_i}.
\]

\medskip
\noindent\underline{Exemple de mulțimi factor}

\noindent\textbf{(a) Numere raționale.}
Considerăm $\mathbb Z\times\mathbb Z$ și relația
\[
 (m,n)\sim(m',n')
 \overset{\mathrm{def}}{\iff}mn'=nm'.
\]
În notițe mulțimea factor este identificată cu
\[
 (\mathbb Z\times\mathbb Z)/{\sim}
   \overset{\mathrm{def}}{=}\mathbb Q,
 \qquad
 r\in\mathbb Q,\quad r=\widehat{(m,n)}.
\]
% Observație: pentru construcția dorită trebuie n,n' != 0.

\noindent\textbf{(b) Numere reale prin șiruri Cauchy.}
Fie
\[
 \mathcal C=\bigl\{(x_n)_{n\geq1}\text{ cu }x_n\in\mathbb Q
       \mid (x_n)_{n\geq1}\text{ este șir Cauchy}\bigr\}.
\]
Definim
\[
 (x_n)_n\sim(y_n)_n
 \overset{\mathrm{def}}{\iff}
 \lim_{n\to\infty}(x_n-y_n)=0.
\]
Aceasta este o relație de echivalență; cu notația manuscrisului,
\[
 \mathcal C/{\sim}\ \overset{\mathrm{not.}}{=}\ \mathbb R,
 \qquad
 r=\widehat{(x_n)_n}\in\mathbb R.
\]

\medskip
\noindent\textbf{(i)} Pe $\mathbb R$ definim
\[
 x\sim y\iff x-y\in\mathbb Z.
\]
Aceasta este relație de echivalență și\enspace\(\mathbb R/{\sim}\cong[0,1)\).
Schița proiecției factor:
\[
\begin{tikzcd}[column sep=large,row sep=small]
 \mathbb R\arrow[r,"p_{\sim}"]\arrow[d,"{x\mapsto x-\lfloor x\rfloor}"']
    & \mathbb R/{\sim} \\
 {[0,1)}\arrow[ur,dashed,"\cong"'] &
\end{tikzcd}
\]
% Eticheta aplicației „parte fracționară” explicitează săgeata manuscrisului.

\noindent\textbf{(ii)} Pe $\mathbb C$ definim
\[
 z_1\sim z_2\overset{\mathrm{def}}{\iff}|z_1|=|z_2|.
\]
Aceasta este relație de echivalență pe $\mathbb C$, iar\enspace\(\mathbb C/{\sim}\cong[0,+\infty)\).

\noindent\textbf{(iii)} Pe $\mathbb C^*$ definim
\[
 z_1\sim z_2\overset{\mathrm{def}}{\iff}
 \arg(z_1)=\arg(z_2).
\]
Aceasta este relație de echivalență pe $\mathbb C^*$ și
\[
 \mathbb C^*/{\sim}\;\cong\;S^1,
 \qquad S^1=\{z\in\mathbb C\mid |z|=1\}.
\]
% Egalitatea argumentelor necesită o alegere consecventă/modulo 2pi.

\noindent\textbf{(iv)} Pe $\mathbb R$ definim
\[
 x\sim y\iff x-y\in\mathbb Q.
\]
Manuscrisul notează $\mathbb R/{\sim}\cong\mathbb R$.

\endgroup
; fără preambul sau pagină nouă.
% NOTE EDITORIALE:
% - Primele rânduri continuă proprietatea de universalitate de pe pagina 4.
%   Termenul din al doilea criteriu este suprascris și poate fi „inj.” ori
%   „bij.”. Am adoptat provizoriu „injectivă”; dacă scria „bijectivă”,
%   ar trebui inclusă și condiția de surjectivitate.
% - În exemplul numerelor raționale apare Z x Z, fără excluderea
%   numitorului zero. Înmulțirea în cruce nu dă o relație de echivalență
%   pe întregul Z x Z; nu am adăugat tacit restricția necesară.
% - Șirurile Cauchy sunt interpretate ca șiruri cu termeni în Q,
%   potrivit notației compacte din manuscris.
% - Schița proiecției din exemplul (i) a fost refăcută ca diagramă.
% - Însemnările foarte palide din marginea dreaptă și cercul hașurat
%   nu au un rol matematic sigur și nu sunt incluse în textul principal.
% - Identificarea R/~ cu R de la (iv) este abstractă, nu canonică.

\subsection*{Seminar 3 - Pagina 5}
\begingroup
\setlength{\abovedisplayskip}{5pt}
\setlength{\belowdisplayskip}{5pt}
\setlength{\abovedisplayshortskip}{3pt}
\setlength{\belowdisplayshortskip}{3pt}

\noindent\textbf{Proprietatea de universalitate (continuare).}
Pentru aplicația $\overline f$ de pe pagina anterioară,
\[
\begin{aligned}
 \overline f\text{ este surjectivă}
   &\iff f\text{ este surjectivă},\\
 \overline f\text{ este injectivă}
   &\iff {\sim}\;=R_f.
\end{aligned}
\]
Mai mult, pentru un sistem de reprezentanți $(x_i)_{i\in I}$,
\[
 M=\bigcup_{i\in I}\widehat{x_i}.
\]

\medskip
\noindent\underline{Exemple de mulțimi factor}

\noindent\textbf{(a) Numere raționale.}
Considerăm $\mathbb Z\times\mathbb Z$ și relația
\[
 (m,n)\sim(m',n')
 \overset{\mathrm{def}}{\iff}mn'=nm'.
\]
În notițe mulțimea factor este identificată cu
\[
 (\mathbb Z\times\mathbb Z)/{\sim}
   \overset{\mathrm{def}}{=}\mathbb Q,
 \qquad
 r\in\mathbb Q,\quad r=\widehat{(m,n)}.
\]
% Observație: pentru construcția dorită trebuie n,n' != 0.

\noindent\textbf{(b) Numere reale prin șiruri Cauchy.}
Fie
\[
 \mathcal C=\bigl\{(x_n)_{n\geq1}\text{ cu }x_n\in\mathbb Q
       \mid (x_n)_{n\geq1}\text{ este șir Cauchy}\bigr\}.
\]
Definim
\[
 (x_n)_n\sim(y_n)_n
 \overset{\mathrm{def}}{\iff}
 \lim_{n\to\infty}(x_n-y_n)=0.
\]
Aceasta este o relație de echivalență; cu notația manuscrisului,
\[
 \mathcal C/{\sim}\ \overset{\mathrm{not.}}{=}\ \mathbb R,
 \qquad
 r=\widehat{(x_n)_n}\in\mathbb R.
\]

\medskip
\noindent\textbf{(i)} Pe $\mathbb R$ definim
\[
 x\sim y\iff x-y\in\mathbb Z.
\]
Aceasta este relație de echivalență și\enspace\(\mathbb R/{\sim}\cong[0,1)\).
Schița proiecției factor:
\[
\begin{tikzcd}[column sep=large,row sep=small]
 \mathbb R\arrow[r,"p_{\sim}"]\arrow[d,"{x\mapsto x-\lfloor x\rfloor}"']
    & \mathbb R/{\sim} \\
 {[0,1)}\arrow[ur,dashed,"\cong"'] &
\end{tikzcd}
\]
% Eticheta aplicației „parte fracționară” explicitează săgeata manuscrisului.

\noindent\textbf{(ii)} Pe $\mathbb C$ definim
\[
 z_1\sim z_2\overset{\mathrm{def}}{\iff}|z_1|=|z_2|.
\]
Aceasta este relație de echivalență pe $\mathbb C$, iar\enspace\(\mathbb C/{\sim}\cong[0,+\infty)\).

\noindent\textbf{(iii)} Pe $\mathbb C^*$ definim
\[
 z_1\sim z_2\overset{\mathrm{def}}{\iff}
 \arg(z_1)=\arg(z_2).
\]
Aceasta este relație de echivalență pe $\mathbb C^*$ și
\[
 \mathbb C^*/{\sim}\;\cong\;S^1,
 \qquad S^1=\{z\in\mathbb C\mid |z|=1\}.
\]
% Egalitatea argumentelor necesită o alegere consecventă/modulo 2pi.

\noindent\textbf{(iv)} Pe $\mathbb R$ definim
\[
 x\sim y\iff x-y\in\mathbb Q.
\]
Manuscrisul notează $\mathbb R/{\sim}\cong\mathbb R$.

\endgroup
; fără preambul sau pagină nouă.
% NOTE EDITORIALE:
% - Primele rânduri continuă proprietatea de universalitate de pe pagina 4.
%   Termenul din al doilea criteriu este suprascris și poate fi „inj.” ori
%   „bij.”. Am adoptat provizoriu „injectivă”; dacă scria „bijectivă”,
%   ar trebui inclusă și condiția de surjectivitate.
% - În exemplul numerelor raționale apare Z x Z, fără excluderea
%   numitorului zero. Înmulțirea în cruce nu dă o relație de echivalență
%   pe întregul Z x Z; nu am adăugat tacit restricția necesară.
% - Șirurile Cauchy sunt interpretate ca șiruri cu termeni în Q,
%   potrivit notației compacte din manuscris.
% - Schița proiecției din exemplul (i) a fost refăcută ca diagramă.
% - Însemnările foarte palide din marginea dreaptă și cercul hașurat
%   nu au un rol matematic sigur și nu sunt incluse în textul principal.
% - Identificarea R/~ cu R de la (iv) este abstractă, nu canonică.

\subsection*{Seminar 3 - Pagina 5}
\begingroup
\setlength{\abovedisplayskip}{5pt}
\setlength{\belowdisplayskip}{5pt}
\setlength{\abovedisplayshortskip}{3pt}
\setlength{\belowdisplayshortskip}{3pt}

\noindent\textbf{Proprietatea de universalitate (continuare).}
Pentru aplicația $\overline f$ de pe pagina anterioară,
\[
\begin{aligned}
 \overline f\text{ este surjectivă}
   &\iff f\text{ este surjectivă},\\
 \overline f\text{ este injectivă}
   &\iff {\sim}\;=R_f.
\end{aligned}
\]
Mai mult, pentru un sistem de reprezentanți $(x_i)_{i\in I}$,
\[
 M=\bigcup_{i\in I}\widehat{x_i}.
\]

\medskip
\noindent\underline{Exemple de mulțimi factor}

\noindent\textbf{(a) Numere raționale.}
Considerăm $\mathbb Z\times\mathbb Z$ și relația
\[
 (m,n)\sim(m',n')
 \overset{\mathrm{def}}{\iff}mn'=nm'.
\]
În notițe mulțimea factor este identificată cu
\[
 (\mathbb Z\times\mathbb Z)/{\sim}
   \overset{\mathrm{def}}{=}\mathbb Q,
 \qquad
 r\in\mathbb Q,\quad r=\widehat{(m,n)}.
\]
% Observație: pentru construcția dorită trebuie n,n' != 0.

\noindent\textbf{(b) Numere reale prin șiruri Cauchy.}
Fie
\[
 \mathcal C=\bigl\{(x_n)_{n\geq1}\text{ cu }x_n\in\mathbb Q
       \mid (x_n)_{n\geq1}\text{ este șir Cauchy}\bigr\}.
\]
Definim
\[
 (x_n)_n\sim(y_n)_n
 \overset{\mathrm{def}}{\iff}
 \lim_{n\to\infty}(x_n-y_n)=0.
\]
Aceasta este o relație de echivalență; cu notația manuscrisului,
\[
 \mathcal C/{\sim}\ \overset{\mathrm{not.}}{=}\ \mathbb R,
 \qquad
 r=\widehat{(x_n)_n}\in\mathbb R.
\]

\medskip
\noindent\textbf{(i)} Pe $\mathbb R$ definim
\[
 x\sim y\iff x-y\in\mathbb Z.
\]
Aceasta este relație de echivalență și\enspace\(\mathbb R/{\sim}\cong[0,1)\).
Schița proiecției factor:
\[
\begin{tikzcd}[column sep=large,row sep=small]
 \mathbb R\arrow[r,"p_{\sim}"]\arrow[d,"{x\mapsto x-\lfloor x\rfloor}"']
    & \mathbb R/{\sim} \\
 {[0,1)}\arrow[ur,dashed,"\cong"'] &
\end{tikzcd}
\]
% Eticheta aplicației „parte fracționară” explicitează săgeata manuscrisului.

\noindent\textbf{(ii)} Pe $\mathbb C$ definim
\[
 z_1\sim z_2\overset{\mathrm{def}}{\iff}|z_1|=|z_2|.
\]
Aceasta este relație de echivalență pe $\mathbb C$, iar\enspace\(\mathbb C/{\sim}\cong[0,+\infty)\).

\noindent\textbf{(iii)} Pe $\mathbb C^*$ definim
\[
 z_1\sim z_2\overset{\mathrm{def}}{\iff}
 \arg(z_1)=\arg(z_2).
\]
Aceasta este relație de echivalență pe $\mathbb C^*$ și
\[
 \mathbb C^*/{\sim}\;\cong\;S^1,
 \qquad S^1=\{z\in\mathbb C\mid |z|=1\}.
\]
% Egalitatea argumentelor necesită o alegere consecventă/modulo 2pi.

\noindent\textbf{(iv)} Pe $\mathbb R$ definim
\[
 x\sim y\iff x-y\in\mathbb Q.
\]
Manuscrisul notează $\mathbb R/{\sim}\cong\mathbb R$.

\endgroup
; fără preambul sau pagină nouă.
% NOTE EDITORIALE:
% - Primele rânduri continuă proprietatea de universalitate de pe pagina 4.
%   Termenul din al doilea criteriu este suprascris și poate fi „inj.” ori
%   „bij.”. Am adoptat provizoriu „injectivă”; dacă scria „bijectivă”,
%   ar trebui inclusă și condiția de surjectivitate.
% - În exemplul numerelor raționale apare Z x Z, fără excluderea
%   numitorului zero. Înmulțirea în cruce nu dă o relație de echivalență
%   pe întregul Z x Z; nu am adăugat tacit restricția necesară.
% - Șirurile Cauchy sunt interpretate ca șiruri cu termeni în Q,
%   potrivit notației compacte din manuscris.
% - Schița proiecției din exemplul (i) a fost refăcută ca diagramă.
% - Însemnările foarte palide din marginea dreaptă și cercul hașurat
%   nu au un rol matematic sigur și nu sunt incluse în textul principal.
% - Identificarea R/~ cu R de la (iv) este abstractă, nu canonică.

\subsection*{Seminar 3 - Pagina 5}
\begingroup
\setlength{\abovedisplayskip}{5pt}
\setlength{\belowdisplayskip}{5pt}
\setlength{\abovedisplayshortskip}{3pt}
\setlength{\belowdisplayshortskip}{3pt}

\noindent\textbf{Proprietatea de universalitate (continuare).}
Pentru aplicația $\overline f$ de pe pagina anterioară,
\[
\begin{aligned}
 \overline f\text{ este surjectivă}
   &\iff f\text{ este surjectivă},\\
 \overline f\text{ este injectivă}
   &\iff {\sim}\;=R_f.
\end{aligned}
\]
Mai mult, pentru un sistem de reprezentanți $(x_i)_{i\in I}$,
\[
 M=\bigcup_{i\in I}\widehat{x_i}.
\]

\medskip
\noindent\underline{Exemple de mulțimi factor}

\noindent\textbf{(a) Numere raționale.}
Considerăm $\mathbb Z\times\mathbb Z$ și relația
\[
 (m,n)\sim(m',n')
 \overset{\mathrm{def}}{\iff}mn'=nm'.
\]
În notițe mulțimea factor este identificată cu
\[
 (\mathbb Z\times\mathbb Z)/{\sim}
   \overset{\mathrm{def}}{=}\mathbb Q,
 \qquad
 r\in\mathbb Q,\quad r=\widehat{(m,n)}.
\]
% Observație: pentru construcția dorită trebuie n,n' != 0.

\noindent\textbf{(b) Numere reale prin șiruri Cauchy.}
Fie
\[
 \mathcal C=\bigl\{(x_n)_{n\geq1}\text{ cu }x_n\in\mathbb Q
       \mid (x_n)_{n\geq1}\text{ este șir Cauchy}\bigr\}.
\]
Definim
\[
 (x_n)_n\sim(y_n)_n
 \overset{\mathrm{def}}{\iff}
 \lim_{n\to\infty}(x_n-y_n)=0.
\]
Aceasta este o relație de echivalență; cu notația manuscrisului,
\[
 \mathcal C/{\sim}\ \overset{\mathrm{not.}}{=}\ \mathbb R,
 \qquad
 r=\widehat{(x_n)_n}\in\mathbb R.
\]

\medskip
\noindent\textbf{(i)} Pe $\mathbb R$ definim
\[
 x\sim y\iff x-y\in\mathbb Z.
\]
Aceasta este relație de echivalență și\enspace\(\mathbb R/{\sim}\cong[0,1)\).
Schița proiecției factor:
\[
\begin{tikzcd}[column sep=large,row sep=small]
 \mathbb R\arrow[r,"p_{\sim}"]\arrow[d,"{x\mapsto x-\lfloor x\rfloor}"']
    & \mathbb R/{\sim} \\
 {[0,1)}\arrow[ur,dashed,"\cong"'] &
\end{tikzcd}
\]
% Eticheta aplicației „parte fracționară” explicitează săgeata manuscrisului.

\noindent\textbf{(ii)} Pe $\mathbb C$ definim
\[
 z_1\sim z_2\overset{\mathrm{def}}{\iff}|z_1|=|z_2|.
\]
Aceasta este relație de echivalență pe $\mathbb C$, iar\enspace\(\mathbb C/{\sim}\cong[0,+\infty)\).

\noindent\textbf{(iii)} Pe $\mathbb C^*$ definim
\[
 z_1\sim z_2\overset{\mathrm{def}}{\iff}
 \arg(z_1)=\arg(z_2).
\]
Aceasta este relație de echivalență pe $\mathbb C^*$ și
\[
 \mathbb C^*/{\sim}\;\cong\;S^1,
 \qquad S^1=\{z\in\mathbb C\mid |z|=1\}.
\]
% Egalitatea argumentelor necesită o alegere consecventă/modulo 2pi.

\noindent\textbf{(iv)} Pe $\mathbb R$ definim
\[
 x\sim y\iff x-y\in\mathbb Q.
\]
Manuscrisul notează $\mathbb R/{\sim}\cong\mathbb R$.

\endgroup
